\documentclass[a4paper, 12pt]{article}

\usepackage{utils}

\renewcommand*{\today}{20 April 2026}

\begin{document}

\hotbox{Analyse 4}{CM 15}{\today}

\section{Intégrales dépendantes d'un paramètre (ou fonction définie par une intégrales)}

\begin{definition}
    Soit $I, J$ deux intervalles de $\R$,
    $$
    \funcdef{f}{I}{\R}{x}{\int_J g(x, t) \, dt}
    $$
    est une fonction définie par une intégrale, ou une intégrale dépendante d'un paramètre.
\end{definition}

\begin{theoreme}{Théorème de la convergence dominée}{}
    Soit $J$ un intervalle de $\R$ et $(f_n)_{n \in \N}$ une suite de fonctions continues par morceaux sur $J$.

    Soient $f$ et $\varphi$ deux fonctions continues par morceaux sur $J$ et $\varphi$ à valeurs positives, telles que
    \begin{itemize}
        \item $\forall t \in J, f_n(t) \xrightarrow[n \to +\infty]{} f(t)$
        \item $\forall t \in J, |f_n(t)| \leq \varphi(t)$
        \item $\varphi \text{ est intégrable sur } J$
    \end{itemize}
    Alors
    $$
    \lim_{n \to +\infty} \int_J f_n(t) \, dt = \int_J \lim_{n \to +\infty} f_n(t) \, dt = \int_J f(t) \, dt
    $$
\end{theoreme}

\begin{demonstration}
    Admis, besoin de l'intégrale de Lebesgue pour le démontrer sur un intervalle non borné.
\end{demonstration}

\subsection{Continuité sous le signe somme}

\begin{theoreme}{Théorème de la continuité d'une fonction définie par une intégrale}{}
    Soit $I, J$ deux intervalles de $\R$ et $\funcdef{f}{J \times I}{\R}{(x, t)}{f(x, t)}$ une fonction dont on suppose que
    \begin{itemize}
        \item $\forall x \in J, f(x, \cdot)$ est continue (?) sur $I$,
        \item $\forall t \in I, f(\cdot, t)$ est continue par morceaux sur $J$,
        \item il existe une fonction $\varphi$ continue par morceaux et intégrable sur $I$ telles que $\forall (x, t) \in J \times I, |f(x, t)| \leq \varphi(t)$.
    \end{itemize}
    Alors la fonction $\funcdef{F}{J}{\R}{x}{\int_I f(x, t) \, dt}$ est continue sur $J$.
\end{theoreme}

\begin{demonstration}
    On fixe $x \in J$, la fonction $f(x, \cdot)$ est continue par morceaux sur $I$ et majorée par $\varphi$ où $\varphi$ est intégrable et continue par morceaux sur $I$.
    Donc $F(x) = \int_I f(x, t) \, dt$ est bien définie.
    $$
    |F(x)| = \left| \int_I f(x, t) \, dt \right| \leq \int_I |f(x, t)| \, dt \leq \int_I \varphi(t) \, dt < +\infty
    $$

    Soit $(x_n)_{n \in \N}$ une suite convergeant vers $x$ dans $J$.
    
    On pose $\forall n \in \N, g_n(t) := f(x_n, t)$, par hypothèse, $\forall t \in I, g_n$ est continue par morceaux sur $I$ et $x \mapsto f(x, t)$ est continue par morceaux sur $J$.

    Donc $f(x_n, t) \xrightarrow[n \to +\infty]{} f(x, t)$ pour tout $t \in I$, ainsi $g_n(t) \xrightarrow[n \to +\infty]{} f(x, t)$ pour tout $t \in I$.

    Donc $g_n$ converge simplement vers $f(x, \cdot)$ sur $I$.

    Et enfin, pour tout $(x, t) \in J \times I$, $|g_n(t)| = |f(x_n, t)| \leq \varphi(t)$, et $\varphi$ est intégrable sur $I$.

    Cela vérifie les hypothèses du théorème de la convergence dominée, et on en déduit que
    $$
    \lim_{n \to +\infty} \int_I g_n(t) \, dt = \int_I f(x, t) \, dt \iff \lim_{n \to +\infty} \int_I f(x_n, t) \, dt = \int_I f(x, t) \, dt \iff \lim_{n \to +\infty} F(x_n) = F(x)
    $$
    Donc $F$ est continue en $x$ et comme $x$ est arbitraire, $F$ est continue par morceaux sur $J$.
\end{demonstration}

\begin{theoreme}{Passage à la limite sous le signe somme}{}
    Soit $I, J$ deux intervalles de $\R$, $a \in \overline{J}$ (l'adhérence de $J$) et $\funcdef{f}{J \times I}{\R}{(x, t)}{f(x, t)}$ une fonction dont on suppose que
    \begin{itemize}
        \item $\forall x \in J, f(x, \cdot)$ est continue sur $I$,
        \item $\forall t \in I, x \mapsto f(x, t)$ a une limite quand $x \to a$,
        \item il existe une fonction $\varphi$ positives et intégrable sur $I$ telles que $\forall (x, t) \in J \times I, |f(x, t)| \leq \varphi(t)$.
    \end{itemize}
    Alors la fonction $\funcdef{F}{J}{\R}{x}{\int_I f(x, t) \, dt}$ a une limite quand $x \to a$ et
    $$
    \lim_{x \to a} F(x) = \int_I \lim_{x \to a} f(x, t) \, dt
    $$
\end{theoreme}

\begin{demonstration}
    On suppose que $a \neq \mp \infty$, sinon, on peut faire un changement de variable pour ramener à ce cas.

    On considère $J \cup \{a\}$, $\forall (x, t) \in (J \cup \{a\}) \times I, g(x, t) := \begin{cases}
    f(x, t) & \text{si } x \in J \\
    \lim_{x \to a} f(x, t) & \text{si } x = a
    \end{cases}$.
    En faisant des trucs et en utilisant le théorème de la continuité d'une fonction définie par une intégrale,
    pour $g$ on a bien le résultat $\funcdef{G}{J \cup \{a\}}{\R}{x}{\int_I g(x, t) \, dt}$ est continue en $a$
    d'où
    $$
    \lim_{x \to a} \int_I f(x, t) \, dt = \int_I \lim_{x \to a} f(x, t) \, dt
    $$
\end{demonstration}

\begin{theoreme}{Théorème de Leibniz}{}
    Soit $I, J$ deux intervalles de $\R$, $\funcdef{f}{J \times I}{\R}{(x, t)}{f(x, t)}$ une fonction dont on suppose que 

    \begin{itemize}
        \item $\forall x \in J, f(x, \cdot)$ est continue par morceaux et intégrable sur $I$
        \item $\forall (x, t) \in J \times I, \frac{\partial f}{\partial x}(x, t)$ existe
        \item $\forall x \in J, t \mapsto \frac{\partial f}{\partial x}(x, t)$ est continue par morceaux sur $I$
        \item $\forall t \in I, x \mapsto \frac{\partial f}{\partial x}(x, t)$ est continue sur $J$
        \item il existe une fonction $\varphi$ définie sur $I$, positive et intégrable sur $I$ telles que $\forall (x, t) \in J \times I, \left| \frac{\partial f}{\partial x}(x, t) \right| \leq \varphi(t)$
    \end{itemize}
    Alors la fonction $\funcdef{F}{J}{\R}{x}{\int_I f(x, t) \, dt}$ est $C^1$ sur $J$ et $\forall x \in J, F'(x) = \int_I \frac{\partial f}{\partial x}(x, t) \, dt$.
\end{theoreme}

\begin{demonstration}
    Comme $\forall x \in J, t \mapsto f(x, t)$ est intégrable sur $I$ alors $\forall x \in J, F(x)$ est bien définie.

    $\forall (x, t) \in J \times I, a \in J, g(x, t) := \begin{cases}
    \frac{f(x, t) - f(a, t)}{x - a} & \text{si } x \neq a \\
    \frac{\partial f}{\partial x}(a, t) & \text{si } x = a
    \end{cases}$
    De sorte que si $x \neq a$, $\frac{F(x) - F(a)}{x - a} = \int_I g(x, t) \, dt$.

    Pour $x \in J$, $t \mapsto g(x, t)$ est continue par morceaux sur $I$.
    Pour $t \in I$, $x \mapsto g(x, t)$ est continue sur $J$ (car $g(x, t)$ est défini comme limite de fonctions continues).
    On a $|g(x, t)| = \left| \frac{f(x, t) - f(a, t)}{x - a} \right| \leq \sup_{y \in J} \left| \frac{\partial f}{\partial x}(y, t) \right| = \varphi_1(t)$, et $\varphi_1$ est intégrable et positive sur $I$.
    
    On peut appliquer le théorème de passage à la limite sous le signe somme, et on en déduit que
    $$
    \lim_{x \to a} \frac{F(x) - F(a)}{x - a} = \lim_{x \to a} \int_I g(x, t) \, dt = \int_I \lim_{x \to a} g(x, t) \, dt = \int_I \frac{\partial f}{\partial x}(a, t) \, dt
    $$
    Donc $F$ est dérivable en $a$ et $F'(a) = \int_I \frac{\partial f}{\partial x}(a, t) \, dt$.
    Puis en appliquant le théorème de la continuité d'une fonction définie par une intégrale, on en déduit que $F'$ est continue sur $J$, et donc $F$ est de classe $C^1$ sur $J$.
\end{demonstration}

\begin{theoreme}{}{}
    Soit $I, J$ deux intervalles de $\R$, $\funcdef{f}{J \times I}{\R}{(x, t)}{f(x, t)}$ une fonction dont on suppose que
    \begin{itemize}
        \item $\forall x \in J, f(x, \cdot)$ est continue par morceaux et intégrable sur $I$
        \item $\forall (x, t) \in J \times I, f \text{ admet des dérivées partielles successives jusqu'à l'ordre } n > 1$
        \item $\forall k \in \llbracket 1, n \rrbracket, \forall x \in J, t \mapsto \frac{\partial^k f}{\partial x^k}(x, t)$ est continue par morceaux sur $I$
        \item $\forall k \in \llbracket 1, n \rrbracket, \forall t \in I, x \mapsto \frac{\partial^k f}{\partial x^k}(x, t)$ est continue sur $J$
        \item $\forall k \in \llbracket 1, n \rrbracket, \forall (x, t) \in J \times I, \left| \frac{\partial^k f}{\partial x^k}(x, t) \right| \leq \varphi_k(t)$ où $\varphi_k$ est une fonction positive et intégrable sur $I$
    \end{itemize}
    Alors la fonction $\funcdef{F}{J}{\R}{x}{\int_I f(x, t) \, dt}$ est de classe $C^n$ sur $J$ et $\forall k \in \llbracket 1, n \rrbracket, \forall x \in J, F^{(k)}(x) = \int_I \frac{\partial^k f}{\partial x^k}(x, t) \, dt$.
\end{theoreme}

\begin{demonstration}
    Par récurrence sur $n$.
\end{demonstration}

\end{document}